\clearpage
\phantomsection
\chapter*{Висновок}
\addcontentsline{toc}{chapter}{Висновок}
\label{chapter:conclusion}

У ході лабораторної роботи виконано повний цикл підготовки корпусу даних для RAG-системи в такій послідовності:
\begin{enumerate}[wide]
    \item завантажено вихідні дані з публічних джерел (статті вікі GTNH, JSON-файли квестів із Git та дамп предметів і рецептів через nesql-exporter) із документуванням їхніх ліцензій;
    \item виконано окреме очищення текстів вікі та квестів від службової, візуальної й технічної розмітки;
    \item збагачено очищені квести єдиними ідентифікаторами та назвами предметів із бази nesql;
    \item реалізовано модульне чанкування статей вікі та квестів за патерном <<Стратегія>> з формуванням структурованого payload (WikiChunk, QuestChunk) для подальшого індексування у Qdrant.
\end{enumerate}

У результаті отримано 9854 очищених та збагачених чанки даних (6020 зі статей вікі та 3834 із квестів), готових до індексування у Qdrant, а також сформовано data card корпусу.